\Needspace{7\baselineskip}
\pdfbookmark[0]{Notices. Schwab and Metternich: Parallel Lines. 1816.}{anzeigen-schwab-metternich-1816}
\begin{center}
\rule{.85\textwidth}{.3pt}\par
Göttingische gelehrte Anzeigen.\quad 20 April 1816.\par
\rule{.85\textwidth}{.3pt}
\end{center}
\hspace*{\parindent}\textit{Commentatio in primum elementorum} \textsc{Euclidis} \textit{librum, qua veritatem geometriae principiis ontologicis niti evincitur, omnesque propositiones, axiomatum geometricorum loco habitae, demonstrantur. Auctore} \textsc{J. C. Schwab}, \textit{Regi Württembergiae a consiliis aulicis secretioribus, academiae scientiarum Petropolitanae, Berolinensis et Harlemensis Sodali.} (65 pages in octavo.) Stuttgart 1814. Typis \textsc{J. F. Steinkopf}.

\textit{Vollständige Theorie der Parallel-Linien. Nebst einem Anhange, in welchem der erste Grundsatz zur Technik der geraden Linie angegeben wird. Herausgegeben von} \textsc{Matthias Metternich}, \textit{Doctor der Philosophie, Professor der Mathematik, Mitglied der gelehrten Gesellschaft nützlicher Wissenschaften zu Erfurt.} 44 pages in octavo. Mainz 1815. Published at the author's expense and distributed by \textsc{Florian Kupferberg}.

Few topics in mathematics have likely been written about as much as the gap at the beginning of geometry in the foundations of the theory of parallel lines. A year seldom goes by without a new attempt to fill this gap. Yet if we are to speak honestly and openly, we cannot say that in substance we have gone any farther than \textsc{Euclid} had %
